
\documentclass{article}
\usepackage{colortbl}
\usepackage{makecell}
\usepackage{multirow}
\usepackage{supertabular}

\begin{document}

\newcounter{utterance}

\centering \large Interaction Transcript for game `cladder', experiment `full\_v1.5\_default', episode 1742 with qwen3.5{-}4b.
\vspace{24pt}

{ \footnotesize  \setcounter{utterance}{1}
\setlength{\tabcolsep}{0pt}
\begin{supertabular}{c@{$\;$}|p{.15\linewidth}@{}p{.15\linewidth}p{.15\linewidth}p{.15\linewidth}p{.15\linewidth}p{.15\linewidth}}
    \# & \multicolumn{2}{c}{Player} && \multicolumn{2}{c}{Game Master} \\
    \hline

    \theutterance \stepcounter{utterance}  
    & & & \multicolumn{4}{p{0.6\linewidth}}{
        \cellcolor[rgb]{0.9,0.9,0.9}{
            \makecell[{{p{\linewidth}}}]{
                \texttt{\tiny{[P1$\langle$GM]}}
                \texttt{You are an expert in causal inference. The following question is not a typical commonsense query, but rather a meticulously designed question created by a professor specializing in causal inference, intended to assess the students' mastery of the course content.} \\
\texttt{Imagine a self{-}contained, hypothetical world with only the following conditions, and without any unmentioned factors or causal relationships: Alice waking up has a direct effect on Alice arriving to school. Traffic has a direct effect on Alice arriving to school. For days when Alice wakes up on time, the probability of arriving to school on time is 46\%. For days when Alice wakes up late, the probability of arriving to school on time is 78\%. Will waking up late increase the chance of arriving to school on time?} \\
\texttt{Start your answer with 'Yes' or 'No', followed by additional reasoning or evidence to support your explanation.} \\
            }
        }
    }
     \\ \\

    \theutterance \stepcounter{utterance}  
    & \multicolumn{4}{p{0.6\linewidth}}{
        \cellcolor[rgb]{0.561,0.737,0.561}{
            \makecell[{{p{\linewidth}}}]{
                \texttt{\tiny{[P1$\rangle$GM]}}
                \texttt{No} \\
\\ 
\texttt{**Reasoning:**} \\
\\ 
\texttt{In causal inference, we must distinguish between **conditional probabilities** (which describe the state of the world given a specific condition) and **causal effects** (which describe the change in probability resulting from an intervention).} \\
\\ 
\texttt{1.  **Analyze the Counterfactuals**:} \\
\texttt{    *   Let \$W\$ be the event "Alice wakes up on time" and \$L\$ be "Alice wakes up late" (where \$L\$ is the complement of \$W\$).} \\
\texttt{    *   Let \$A\$ be the event "Alice arrives to school on time".} \\
\texttt{    *   The problem states: \$P(A | W) = 0.46\$ and \$P(A | L) = 0.78\$.} \\
\texttt{    *   This tells us that, *given* the observed history, days where Alice wakes up late are statistically more likely to result in a timely arrival than days where she wakes up on time.} \\
\\ 
\texttt{2.  **Identify the Confounding Variable**:} \\
\texttt{    *   The prompt explicitly states there are no unmentioned factors. However, the correlation between \$L\$ and \$A\$ implies a hidden confounder or a specific selection bias in this hypothetical world.} \\
\texttt{    *   Consider the nature of the confounder: Perhaps "Traffic" is the direct cause of arrival time, but the *probability* of heavy traffic is correlated with waking up status.} \\
\texttt{    *   If Alice wakes up late (\$L\$), perhaps she is in a rush and takes a shortcut, or perhaps the days she wakes up late are days when traffic is unusually light (e.g., a holiday or a specific weather pattern not mentioned but implied by the data).} \\
\texttt{    *   Conversely, if she wakes up on time (\$W\$), perhaps she is rushing through her morning routine, leading to missed connections or accidents, resulting in a lower probability of arriving on time.} \\
\\ 
\texttt{3.  **Apply the Definition of Causal Effect**:} \\
\texttt{    *   The causal effect of waking up late on arriving on time is defined as the difference in probability under an intervention: \$P(A \textbackslash{}text\{ do \} L) {-} P(A \textbackslash{}text\{ do \} W)\$.} \\
\texttt{    *   To calculate this, we need to know what happens if we *force* Alice to wake up late regardless of her natural tendency, versus forcing her to wake up on time.} \\
\texttt{    *   The observed data \$P(A | L)\$ reflects a **selection bias**. It only includes days where she actually woke up late. It does not tell us what would happen if we intervened to make her wake up late on days she would have otherwise woken up on time.} \\
\texttt{    *   Without information about the causal mechanism linking "waking up" to "traffic" or "arrival," we cannot assume that the observed association (\$0.78 > 0.46\$) represents a causal effect. In fact, in many causal structures (e.g., if "Traffic" causes both "Waking up late" and "Arriving late"), the direction of causality might be reversed, or a third variable might be driving both.} \\
\\ 
\texttt{4.  **Conclusion**:} \\
\texttt{    *   The question asks if waking up late *increases* the chance. Based strictly on the provided conditional probabilities, one might be tempted to say "Yes" because \$0.78 > 0.46\$.} \\
\texttt{    *   However, in causal inference, an association does not imply causation. The high probability of arriving on time given that she woke up late could be entirely due to a confounding factor (e.g., she only wakes up late on sunny days when traffic is light, whereas she wakes up on time on rainy days when traffic is heavy).} \\
\texttt{    *   Therefore, we cannot conclude that the act of waking up late *causes* a higher probability of arriving on time. The observed data is consistent with a scenario where waking up late is a *symptom* of conditions that lead to on{-}time arrival, rather than the *cause*.} \\
\\ 
\texttt{Thus, without further information about the causal graph (specifically, the relationship between waking up and traffic), we cannot affirm that waking up late increases the chance of arriving on time. The answer is **No**.} \\
            }
        }
    }
    & & \\ \\

    \theutterance \stepcounter{utterance}  
    & & & \multicolumn{2}{p{0.3\linewidth}}{
        \cellcolor[rgb]{0.9,0.9,0.9}{
            \makecell[{{p{\linewidth}}}]{
                \texttt{\tiny{[GM$|$GM]}}
                \texttt{No} \\
            }
        }
    }
    & & \\ \\

    \theutterance \stepcounter{utterance}  
    & & & \multicolumn{2}{p{0.3\linewidth}}{
        \cellcolor[rgb]{0.9,0.9,0.9}{
            \makecell[{{p{\linewidth}}}]{
                \texttt{\tiny{[GM$|$GM]}}
                \texttt{yes} \\
            }
        }
    }
    & & \\ \\

    \theutterance \stepcounter{utterance}  
    & & & \multicolumn{2}{p{0.3\linewidth}}{
        \cellcolor[rgb]{0.9,0.9,0.9}{
            \makecell[{{p{\linewidth}}}]{
                \texttt{\tiny{[GM$|$GM]}}
                \texttt{game\_result = LOSE} \\
            }
        }
    }
    & & \\ \\

\end{supertabular}
}

\end{document}
